\section{Metoda pomiarów}
\label{sec:method}

Praca zawiera 5 pomiarów. Pomiary P1 i P2 są obliczeniami na modelu formalnym. Pomiary P3
i P4 korzystają ze zbioru pozycji testowych. Pomiar P5 korzysta z korpusu tekstów.

\subsection{Zbiór pozycji testowych}
\label{sec:items}

Zbiór ma \pkItems{} pozycji (plik \texttt{data/items.csv}, dodatek \ref{app:items}).
Pozycja zawiera 6 części: zdanie zwykłe, opis 2 odczytań, zdanie PK dla odczytania 1,
zdanie PK dla odczytania 2 oraz 3 sondy. Sonda to sytuacja z pytaniem, na które
odpowiedzią jest \pkq{Tak} albo \pkq{Nie}. Pierwsza sonda leży w strefie szarej: odczytanie
1 daje jedną odpowiedź, a odczytanie 2 daje odpowiedź przeciwną. Sondy druga i trzecia są
kotwicami: oba odczytania dają odpowiedź \pkq{Tak} w kotwicy drugiej i odpowiedź \pkq{Nie}
w kotwicy trzeciej. Kotwice sprawdzają, czy model rozumie zadanie.

Zdanie zwykłe w \pkItemsStd{} pozycjach pochodzi z przykładów standardu PK 0.1. Standard
powstał przed tą pracą, więc autor nie dobierał tych zdań pod wynik pomiaru. Pozostałe
\pkItemsAuthor{} pozycji napisał autor. Zbiór ma 3 rodzaje pozycji.
\begin{itemize}[leftmargin=1.6em]
  \item \textbf{Przeredagowanie} (\pkItemsRewrite{} pozycji): zdanie PK jest innym ciągiem
  słów niż zdanie zwykłe.
  \item \textbf{Konwencja} (\pkItemsConv{} pozycje, reguły B2 i C1): zdanie PK dla odczytania
  1 jest tym samym ciągiem słów co zdanie zwykłe. Odczytanie ustala konwencja, więc
  pozycja ma dodatkowy wariant z legendą.
  \item \textbf{Luka} (\pkItemsGap{} pozycje): zdanie zwykłe spełnia numerowane reguły PK
  i ma 2 odczytania (uwaga \ref{rem:gap}).
\end{itemize}
Każde zdanie PK w zbiorze przechodzi sprawdzarkę z sekcji \ref{sec:check}. Zdanie PK
o największej liczbie słów ma \pkLintPkLongest{} słów.

\subsection{Pomiary P1 i P2}

Pomiar P1 liczy odczytania dokładnie. Funkcja \texttt{count\_parses} liczy drzewa rozbioru
zdania z $k$ frazami przyimkowymi. Funkcja \texttt{bracketings} wylicza nawiasowania ciągu
członów połączonych na przemian słowami \pkq{i} oraz \pkq{lub} i liczy różne funkcje
logiczne. Funkcja \texttt{scope\_orders} liczy różne porządki zasięgu kwantyfikatorów
\pkq{każdy} i \pkq{pewien} ustawionych na przemian. Dla co najwyżej 4 kwantyfikatorów test
sprawdza wynik semantycznie: wylicza każdą relację na dziedzinie o 2 elementach.

Pomiar P2 uruchamia parser Earleya na zdaniach o budowie: rzeczownik, czasownik,
rzeczownik oraz $k$ fraz przyimkowych, dla $k$ od 0 do 40 włącznie. Zdanie ma wtedy
od 3 do \pkParseWordsMax{} słów włącznie. Operacją jest 1 próba wpisania stanu do tablicy
parsera. Parser nie zawiera poprawki Leo \cite{leo1991}. Wykładnik wzrostu pochodzi
z regresji liniowej logarytmu liczby operacji względem logarytmu liczby słów, dla zdań,
które mają co najmniej 21 słów.

\subsection{Pomiar P3: długość}

Pomiar P3 liczy długość zdania zwykłego i 2 zdań PK w 6 miarach: słowa, znaki oraz tokeny
\pkTokTokenizers{} tokenizerów. Tokenizery \texttt{o200k\_base} i \texttt{cl100k\_base}
pochodzą z biblioteki \texttt{tiktoken}; tokenizery modeli \texttt{Llama 3.2} i \texttt{Gemma 3} pochodzą
z plików modeli. Tokenizery modeli Bielik wymagają zgody licencyjnej w serwisie
Hugging Face, więc pomiar pomija tokenizery modeli Bielik. Narzut pozycji to iloraz: średnia długość
2 zdań PK podzielona przez długość zdania zwykłego, minus 1. Pomiar obejmuje \pkTokItems{} pozycji:
przeredagowania i luki. Pozycje typu konwencja mają narzut 0 z definicji; kosztem
konwencji jest legenda. Dla mediany narzutu skrypt liczy przedział ufności 95 \% metodą
bootstrap z 10\,000 losowań.

\subsection{Pomiar P4: sonda strefy szarej}
\label{sec:probemethod}

Model językowy dostaje instrukcję systemową: \pkq{Odpowiadasz na pytania o znaczenie
tekstu. Odpowiedz jednym słowem: Tak albo Nie.} Wiadomość zawiera tekst, sondę i powtórzone
polecenie odpowiedzi 1 słowem. Tekstem jest zdanie zwykłe, zdanie PK dla odczytania 1 albo
zdanie PK dla odczytania 2. Pozycja typu konwencja ma także 2 warianty z legendą. Razem
pomiar ma \pkProbes{} sond, czyli \pkProbesUnique{} różnych zapytań na model.

Pomiar ma 2 wykonane tory.
\begin{itemize}[leftmargin=1.6em]
  \item \textbf{Tor lokalny.} Program \texttt{ollama} uruchamia model na procesorze
  komputera autora, z temperaturą 0. Skrypt odczytuje prawdopodobieństwa 20 najbardziej
  prawdopodobnych tokenów na pierwszej pozycji odpowiedzi. Wielkość $P(\text{Tak})$ to
  masa tokenów słowa \pkq{Tak} podzielona przez sumę mas tokenów słów \pkq{Tak}
  i \pkq{Nie}. Tor lokalny obejmuje każdą z \pkProbesUnique{} sond.
  \item \textbf{Tor Claude.} Polecenie \texttt{claude -p} wywołuje modele Haiku i Sonnet
  bez narzędzi, bez plików reguł autora i bez pamięci sesji; mediana zużycia to
  \pkClaudeTokMedian{} tokenów wejściowych na wywołanie. Interfejs nie podaje prawdopodobieństw tokenów. Skrypt zadaje
  więc każdą sondę 2 razy i liczy udział odpowiedzi \pkq{Tak}. Tor Claude obejmuje tylko
  sondy strefy szarej: \pkProbesGrayUnique{} zapytań i \pkClaudeCalls{} wywołania na model.
\end{itemize}
Plan przewidywał tor trzeci, przez serwis OpenRouter. Tor trzeci nie powstał. Zabezpieczenie
komputera autora zablokowało odczyt klucza dostępu, a autor zostawił blokadę na czas
pomiaru.

Niech $q$ oznacza prawdopodobieństwo odczytania 1 w strefie szarej. Skrypt liczy 5 miar.
\begin{itemize}[leftmargin=1.6em]
  \item \textbf{Niepewność} $H_2(q)=-q\log_2q-(1-q)\log_2(1-q)$, w bitach, dla zdania
  zwykłego i średnio dla 2 zdań PK. Miara wymaga prawdopodobieństw, więc dotyczy toru
  lokalnego.
  \item \textbf{Sterowalność}: różnica $q$ dla zdania PK odczytania 1 i $q$ dla zdania PK
  odczytania 2. Wartość 1 oznacza pełne sterowanie; wartość 0 oznacza brak wpływu.
  \item \textbf{Trafność PK}: udział zdań PK, dla których model wybiera odczytanie
  zamierzone, czyli $q>0{,}5$ dla odczytania 1 i $q<0{,}5$ dla odczytania 2. Poziom
  przypadku wynosi 50 \%.
  \item \textbf{Trafność kotwic}: udział kotwic z odpowiedzią zgodną z kluczem.
  \item \textbf{Rozbieżność czytelników}: udział pozycji, w których modele wybierają różne
  odczytania zdania zwykłego.
\end{itemize}
Dla niepewności skrypt wykonuje dokładny test znaków na pozycjach typu przeredagowanie.
Test nie ma poprawki na liczbę modeli.

\subsection{Pomiar P5: korpus}

Korpus zawiera \pkCorpusActs{} aktów prawnych i 3 teksty pisane w PK. Akty prawne pochodzą
z serwisu \texttt{api.sejm.gov.pl}: Konstytucja RP oraz teksty jednolite Kodeksu cywilnego,
Kodeksu pracy, Prawa o ruchu drogowym i \pkq{Zasad techniki prawodawczej}. Akty prawne
nie są przedmiotem prawa autorskiego; repozytorium zapisuje adresy, sumy SHA-256 i wyniki
zliczeń. Teksty PK to: standard PK 0.1, praca \pkq{Rachunek generatorów} tego samego autora
oraz niniejsza praca.

Skrypt usuwa z aktu nagłówki stron i numery jednostek redakcyjnych. Sprawdzarka liczy
trafienia każdej z \pkLintChecks{} kontroli; miarą jest liczba trafień na 1000 słów.
Trafienie wskazuje konstrukcję objętą regułą PK. Trafienie nie dowodzi, że zdanie ma
2 odczytania w kontekście aktu. Autor ocenił ręcznie losową próbę \pkPrecN{} trafień
z aktów prawnych; próba nie obejmuje kontroli długości zdania. Ocenę wykonał 1 oceniający.

\subsection{Powtarzalność}

Polecenie \texttt{make all} uruchamia testy, pomiary P1, P2, P3 i P5, rysunki i budowę
pliku PDF. Surowe odpowiedzi modeli leżą w plikach \texttt{results/probe\_*.jsonl};
skrypt \texttt{scripts/measure.py} przelicza z tych plików każdą liczbę pomiaru P4. Ciągła
integracja nie wywołuje modeli. Wagi modeli lokalnych identyfikuje nazwa z rejestru
programu \texttt{ollama}; modele Claude nie są dostępne bez konta, więc tor Claude
odtwarza tylko posiadacz konta.

\section{Wyniki}
\label{sec:results}

\subsection{P1: liczba odczytań}

Tabela \ref{tab:counts} i rysunek \ref{fig:growth} podają liczbę odczytań dla 3 konstrukcji.
Zdanie z 8 frazami przyimkowymi ma \pkAttachEight{} odczytania. Ciąg 11 członów połączonych
na przemian słowami \pkq{i} oraz \pkq{lub} ma \pkBracketElevenAll{} nawiasowań, które dają
\pkBracketEleven{} różnych funkcji logicznych. Zdanie z 4 kwantyfikatorami ma
\pkScopeFourAll{} porządki zasięgu, w tym \pkScopeFour{} logicznie różnych. Wzrost jest
wykładniczy w każdej z 3 konstrukcji. Zdanie PK ma 1 odczytanie przy każdej liczbie
miejsc wyboru.

\begin{table}[tp]
\centering
\caption{P1: liczba logicznie różnych odczytań zdania. Kolumna 2: $k$ fraz przyimkowych po
dopełnieniu. Kolumna 3: $k$ spójników \pkq{i} oraz \pkq{lub} na przemian. Kolumna 4:
$k+1$ kwantyfikatorów \pkq{każdy} i \pkq{pewien} na przemian.}
\label{tab:counts}
\begin{tabular}{rrrr}
\toprule
$k$ & frazy przyimkowe & spójniki & kwantyfikatory\\
\midrule
\wiersze{data/table_counts.tex}
\bottomrule
\end{tabular}
\end{table}

\rysunek{growth}{Kombinatoryka: wzrost liczby odczytań z liczbą miejsc wyboru.}{fig:growth}

\subsection{P2: operacje parsera}
\label{sec:resparse}

Rysunek \ref{fig:parser} pokazuje liczbę operacji parsera. Dla zdania o 21 słowach parser
wykonuje \pkParseOpsNatTwenty{} operacji z gramatyką niejednoznaczną i \pkParseOpsPkTwenty{}
operacji z gramatyką PK. To samo zdanie ma \pkParseTreesTwenty{} drzew rozbioru
w gramatyce niejednoznacznej i 1 drzewo w gramatyce PK. Dla zdania
o \pkParseWordsMax{} słowach liczby operacji wynoszą \pkParseOpsNatMax{}
oraz \pkParseOpsPkMax{}; iloraz to \pkParseRatioMax.

Dopasowany wykładnik wzrostu wynosi \pkParseExpNat{} dla gramatyki niejednoznacznej
i \pkParseExpPk{} dla gramatyki PK. Oba wykładniki są mniejsze niż granice teoretyczne
3 i 2, bo zdania pomiaru mają co najwyżej \pkParseWordsMax{} słów. Gramatyka PK jest
jednoznaczna i ma rekursję prawostronną. Parser Earleya bez poprawki Leo przetwarza taką
gramatykę w czasie kwadratowym; z poprawką Leo czas jest liniowy \cite{leo1991}. Pomiar
potwierdza kierunek zmiany i nie potwierdza czasu liniowego.

W zakresie reguły B8 przewaga gramatyki PK w liczbie operacji jest umiarkowana: dla zdania
o 25 słowach iloraz wynosi \pkParseRatioTwentyFive. Główną korzyścią w tym zakresie jest
liczba drzew: 1 drzewo zamiast \pkParseTreesTwentyFive{} drzew.

\rysunek{parser}{Teoria złożoności: zmierzona liczba operacji parsera Earleya dla
gramatyki niejednoznacznej i dla gramatyki PK.}{fig:parser}

\subsection{P3: narzut długości}
\label{sec:restokens}

PK wydłuża tekst. Tabela \ref{tab:tokens} i rysunek \ref{fig:tokens} podają narzut
w 6 miarach. W tokenach \pkTokUnit{} mediana narzutu wynosi \pkTokMedian{} \%, a kwartyle
wynoszą \pkTokQone{} \% i \pkTokQthree{} \%. Przedział ufności 95 \% dla mediany to
od \pkTokCiLo{} \% do \pkTokCiHi{} \% włącznie. Średnia narzutu wynosi \pkTokMean{} \%.
Zdanie PK jest dłuższe niż zdanie zwykłe w \pkTokShareLonger{} \% przypadków. Wynik jest
zgodny w 6 miarach: mediany leżą w przedziale od \pkTokMedMin{} \% do \pkTokMedMax{} \% włącznie.

Najwyższy narzut mają pozycje o terminach liczonych w dniach (reguła D4) oraz pozycja
o mieszaniu spójników (reguła C2). Zdanie PK podaje tam dane, których zdanie zwykłe
nie zawiera: dzień numer 1 albo strukturę listy. Narzut bliski 0 mają pozycje, w których
PK zamienia 1 słowo, na przykład słowo \pkq{wszystkie} z przeczeniem na słowo \pkq{żaden}.

Pozycje typu konwencja nie zmieniają zdania. Kosztem jest legenda, którą autor podaje
1 raz w tekście; legendy zbioru mają co najwyżej \pkLegendTokMax{} tokenów.

\begin{table}[tp]
\centering
\caption{P3: narzut długości PK w procentach długości zdania zwykłego; \pkTokItems{} pozycji.
Kolumna końcowa: udział zdań PK dłuższych niż zdanie zwykłe.}
\label{tab:tokens}
\begin{tabular}{lrrrrrr}
\toprule
Miara & mediana & kwartyl 1 & kwartyl 3 & minimum & maksimum & dłuższe (\%)\\
\midrule
\wiersze{data/table_tokens.tex}
\bottomrule
\end{tabular}
\end{table}

\rysunek{tokens}{Lingwistyka korpusowa: długość zdania zwykłego i zdania PK oraz rozkład
narzutu w 6 miarach.}{fig:tokens}
